\chapter{Introduction}
\label{sec:introduction}

Invasive alien plants alter ecosystems, reduce native biodiversity and
are costly to control once they have spread over large areas.
\emph{Acacia dealbata} (silver wattle), native to south-eastern
Australia, is one such species in Portugal, where several \emph{Acacia}
species are classified as invasive under national law (Decree-Law
No.~92/2019). It forms dense stands that outcompete native vegetation,
and knowing where it currently grows is a prerequisite for any control
effort. Producing an up-to-date map over a large area is nonetheless
difficult: the species must be told apart from visually similar trees
and shrubs in a satellite image, and the land it occupies changes from
year to year.

Satellites make it possible to observe vegetation repeatedly and at
scale. The Sentinel-2 mission, for example, revisits the same location
every few days and records several spectral bands beyond what the
human eye sees, at a resolution of 10 metres per pixel. Mapping a
species from these images means classifying each pixel as that species
or not, using whatever signal best separates it from its surroundings;
for \emph{A.\ dealbata}, this is mainly its distinctive winter
flowering period \cite{domingo2023,martins2016}.

Classifying pixels requires reference data: locations where the
species is known to be present (and, ideally, known to be absent),
used to train and evaluate the classifier. Collecting this data in the
field is slow and expensive, which limits how much of it is ever
available for a given study area -- this is the starting constraint
that motivates this thesis.

Two developments offer a way to work within that constraint. The
first is \emph{Earth Observation foundation models}: models
pre-trained on large volumes of satellite imagery that turn each pixel
into a fixed-length vector, or \emph{embedding}, capturing structure
learned across many places and times, rather than a hand-engineered
spectral feature. TESSERA \cite{feng2025} and AlphaEarth
\cite{brown2025} are two recent examples. For a reader familiar with
machine learning but not Earth Observation, the closest analogy is a
pre-trained language or vision model: instead of engineering features
by hand, a light classifier is trained on top of embeddings produced
by a model that has already learned general-purpose structure from far
more data than any single study could label. The second is
\emph{citizen science}: platforms such as iNaturalist let members of
the public record species observations with a photograph, a date and a
location \cite{gillespie2024}. These records are far cheaper to obtain
than a field survey, but they are opportunistic rather than
systematic: a citizen-science dataset is a set of points where someone
happened to look and recognise the species, not a designed survey, and
it almost never records a confirmed absence.

This thesis asks whether foundation-model embeddings make it possible
to map an invasive plant accurately from this kind of cheap, imperfect
reference data, and how much performance is still gained by adding
progressively costlier, expert-collected data on top of it.

\section{Motivation}
\label{sec:introduction-motivation}

Combining citizen science with Earth Observation foundation models is
attractive because the two weaknesses are complementary in principle:
citizen science is cheap but biased and incomplete, while foundation
models are expensive to produce but, once available, are reused for
free across tasks. Whether that combination actually works for a
specific invasive species in a specific region is an open empirical
question, not something that follows automatically from either
technology on its own.

Several concrete problems make it hard, each already observed in this
thesis's own exploratory analysis of the Coimbra-district iNaturalist
data, and each already associated with findings elsewhere in the
literature:

\begin{table}[h]
\centering
\begin{tabularx}{\textwidth}{|>{\raggedright\arraybackslash}X|>{\raggedright\arraybackslash}X|>{\raggedright\arraybackslash}X|}
\hline
\textbf{Problem} & \textbf{In machine-learning terms} & \textbf{Addressed in} \\
\hline
Observations cluster near roads and trails \cite{dimson2023} &
Sampling / dataset bias & RQ3 \\
\hline
A single observer made most observations (83\% overall, 88\% of the
cleaned presences used in this thesis) & Low training-data diversity &
RQ3; observer balancing \\
\hline
Citizen science records only presences, never confirmed absences
\cite{phillips2009} & Positive-unlabelled learning & Background pool
built from other-species observations (Chapter~\ref{sec:methodology}) \\
\hline
Imprecise locations and mixed pixels \cite{ball2026} & Label noise &
Reference-data pipeline (Chapter~\ref{sec:methodology}) \\
\hline
Limited volume of reference data & Data efficiency & RQ2 \\
\hline
Observations recorded in years other than the satellite imagery
\cite{ball2026} & Temporal mismatch between data and imagery & Single
reference year; disturbance filter \\
\hline
Nearby observations inflate accuracy estimates \cite{karasiak2022} &
Data leakage / spatial autocorrelation & Spatial validation \\
\hline
Weaker transfer of embeddings to new regions \cite{zhu2026} &
Distribution shift & RQ3 \\
\hline
\end{tabularx}
\caption{Practical problems in using citizen-science data with Earth
Observation foundation models, and where each is addressed in this
thesis.}
\label{tab:introduction-problems}
\end{table}

\section{Research Questions}
\label{sec:introduction-rqs}

Each ingredient above -- foundation models for mapping, citizen
science as reference data, and invasive species as the target -- has
been studied on its own or in pairs, but their three-way combination,
together with an explicit evaluation of data sparsity and spatial
bias, has not (Chapter~\ref{sec:related-work}). This thesis's main
question is:

\begin{quote}
How robust and data-efficient are Earth Observation foundation models
for invasive plant mapping when trained with citizen-science reference
data?
\end{quote}

This is treated as a question of degree, not a yes/no answer: EO
foundation models may give more data-efficient representations, but
their advantage may be limited by the quality, quantity and spatial
distribution of citizen-science data. The thesis asks under what
conditions they are useful, not only whether they are better.

Three secondary questions operationalise this, in order of priority --
RQ1 is the prerequisite and core deliverable, and each later question
builds on it and on the same pipeline:

\textbf{RQ1 -- Foundation models vs.\ conventional approaches.} To
what extent do Earth Observation foundation-model representations
improve invasive plant mapping over conventional spectral-temporal
satellite features when trained on citizen-science labels, and how
much performance is gained by progressively adding costlier expert
reference data on top of it?

\textbf{RQ2 -- Data efficiency.} How does the performance of EO
foundation models compare with conventional approaches as
citizen-science reference data are reduced by a filter that keeps the
most recent and most reliable observations?

\textbf{RQ3 -- Spatial sampling bias.} Do EO foundation-model
representations amplify or mitigate the effect of spatial sampling
bias in citizen-science observations, compared with conventional
satellite features, when evaluated across geographically distinct
regions?

RQ1 alone already constitutes a complete contribution, answering the
first and third objectives, and part of the second objective, of the
MAP.INVADER project proposal this thesis is part of; RQ2 reuses the
same pipeline, model and test set at a low extra cost; RQ3 is the most
demanding, since it requires a second, geographically distinct study
region, and is pursued only as time and data allow. The full design
for each question, including how expert reference data enters RQ1 as a
progressively added increment rather than a fixed point of comparison,
is given in Chapter~\ref{sec:methodology}.

\section{Contributions}
\label{sec:introduction-contributions}

Two contributions are intended: one already settled as a methodological
design, and one still pending the experiments that will establish it.

\textbf{A documented, evaluated pipeline} that turns raw
citizen-science point observations into labelled satellite training
data, addressing the proposal's first objective. Each cleaning and
filtering step -- removing obscured coordinates, keeping only
identification-confirmed (\emph{research-grade}) records, discarding
imprecise locations, collapsing repeated observations onto the same
pixel, among others -- is documented individually, and its effect on
the final dataset is measured through a cumulative ablation rather
than assumed. One step in particular, matching what citizen scientists
tend to photograph (often low, understory plants) against what the
satellite actually sees (the top of the canopy), has no precedent
found in the literature reviewed so far
(Chapter~\ref{sec:related-work}) and is treated as a candidate
methodological contribution in its own right.

\textbf{An empirical answer to RQ1--RQ3}, established through the
experiments described in Chapter~\ref{sec:methodology} and reported in
Chapter~\ref{sec:results}: how far conventional and foundation-model
representations differ under citizen-science labels, how much of that
gap closes as costlier expert data is progressively added, how
performance degrades as citizen-science data are reduced, and whether
foundation-model representations amplify or dampen the effect of
spatial sampling bias. No result is anticipated here; what follows in
this thesis is the account of how these questions were answered, not a
restatement of answers assumed in advance.

\section{Thesis Structure}
\label{sec:introduction-structure}
\todo[inline]{To be written last, once Chapters~\ref{sec:related-work}
-- \ref{sec:conclusion} are no longer stubs, so each paragraph previews
content that actually exists rather than a plan.}
